\section{The ratio test, forest updates, and termination}\label{sec:pivots}

\subsection{The largest admissible step}
For brevity let $\vect{z}=(\vect{y},\vect{s})$ denote the standard-form variable vector in this
section only, and let $d$ include vertex and slack components from
\eqref{eq:direction}. With entering variable $q$, every other nonbasic
variable stays zero. Nonnegativity of the basic variables gives
\begin{equation}\label{eq:step}
 t^\star=\min_{b\in\Bset:\ d_b<0}\frac{z_b}{-d_b}.
\end{equation}
Choose a leaving variable attaining this minimum, and set
\begin{equation}\label{eq:pivot}
 z'=z+t^\star d,\qquad
 \Nset'=(\Nset\setminus\{q\})\cup\{b\}.
\end{equation}
The set in \eqref{eq:step} cannot be empty: the full normalized polytope,
including its slacks, is bounded, whereas $d_q=1$ would otherwise yield
an infinite feasible ray. The new objective is
$\rho'=\rho+t^\star\bar c_q$.

If a blocking basic variable is already zero, $t^\star=0$. Such a pivot
changes the basis and the forest but not the point. It can be necessary
even when the current point already has the optimal objective value.

\subsection{All four structural cases}
Forest edges encode \emph{nonbasic} slacks. Hence an entering slack is
removed from the forest, whereas a leaving slack is added to it. Roots
obey the corresponding rule for vertex variables.
\begin{table}[htbp]
\centering\small
\begin{tabular}{llll}
\toprule
Entering & Leaving & Edge update & Root update\\
\midrule
$y_q$ & $y_b$ & $\Fset' =\Fset$ & $\Qset'=(\Qset\setminus\{q\})\cup\{b\}$\\
$y_q$ & $s_f$ & $\Fset'=\Fset\cup\{f\}$ & $\Qset'=\Qset\setminus\{q\}$\\
$s_e$ & $y_b$ & $\Fset'=\Fset\setminus\{e\}$ & $\Qset'=\Qset\cup\{b\}$\\
$s_e$ & $s_f$ & $\Fset'=(\Fset\setminus\{e\})\cup\{f\}$ & $\Qset'=\Qset$\\
\bottomrule
\end{tabular}
\caption{Updates determined by the entering and leaving variable types.}
\label{tab:updates}
\end{table}
Removing a forest edge splits a tree; adding a leaving slack joins the
appropriate components after any split. A root exchange can transfer
which component is positive. One should not guess these changes from an
improvement in the objective: the ratio test determines the valid exchange.

\begin{theorem}[Correctness of one pivot]\label{thm:pivot}
Starting from a feasible forest basis, a positive reduced-cost entry and
the full ratio test \eqref{eq:step} produce a feasible basis represented
by \cref{tab:updates}. The objective does not decrease.
\end{theorem}
\begin{proof}
\Cref{prop:direction} identifies exactly the ordinary simplex direction.
The ratio test preserves nonnegativity of every basic variable; equality
constraints and the remaining nonbasic zero values are preserved by the
direction. The pivot element is nonzero because $d_b<0$, so exchanging
the columns yields a nonsingular basis. Its nonbasic complement is exactly
\eqref{eq:pivot}, hence \cref{thm:basis} proves the forest and root
properties without additional assumptions on orientation. The objective
change is $t^\star\bar c_q\ge0$.
\end{proof}

Rebuilding the new components and evaluating \eqref{eq:basic-solution}
must reproduce $z'$. This is a particularly useful implementation check.
A separate heuristic ``tree feasibility'' test cannot replace the complete
ratio test, which includes all decreasing vertex variables and all decreasing
basic slacks.

\subsection{A finite reference algorithm}
Fix a total ordering of all vertex and slack variables. For the reference
version use Bland's rule: the smallest-index nonbasic variable with
$\bar c_q>0$ enters; among all minimum-ratio blockers the smallest-index
basic variable leaves. Both parts of this rule matter for the standard
finiteness guarantee \citep{Bland1977}.

\begin{algorithm}[htbp]
\caption{\FSC{} ratio oracle, exact-arithmetic reference}
\label{alg:forest}
\begin{algorithmic}[1]
\Require Nonempty DAG, profits $p$, positive weights $w$, fixed variable order
\Ensure Optimal ratio $\rho$, an optimal connected closure $T_*$, dual certificate
\State Choose a sink $v$; set $\Fset\gets\emptyset$, $\Qset\gets\Vset\setminus\{v\}$
\Loop
  \State Recover trees, roots and $T_*$; root the trees
  \State Compute subtree $p,w$ sums, basic values, and $\rho$ from \eqref{eq:basic-solution}
  \State Price all nonbasic variables using \eqref{eq:rc-node}--\eqref{eq:rc-edge}
  \If{all reduced costs are nonpositive}
    \State \Return $\rho,T_*$ and the certificate \eqref{eq:forest-dual}
  \EndIf
  \State Select the smallest-index improving nonbasic variable $q$
  \State Compute $h$, vertex directions and all slack directions by \eqref{eq:direction}
  \State Compute $t^\star$ and the smallest-index tied leaving variable $b$ by \eqref{eq:step}
  \State Update $\Fset,\Qset$ using \cref{tab:updates}
\EndLoop
\end{algorithmic}
\end{algorithm}

\begin{theorem}[Finite exact-arithmetic oracle]\label{thm:finite}
\Cref{alg:forest} terminates at an optimum of \eqref{eq:ratio}.
\end{theorem}
\begin{proof}
Initialization is feasible by \eqref{eq:initial}. Each iteration is the
ordinary primal simplex algorithm with Bland's entering and leaving rules,
by \cref{prop:direction,thm:pivot}. The normalized polytope is nonempty
and bounded. Bland's finite pivoting theorem therefore applies. At
termination \cref{prop:certificate} provides an optimality certificate,
and \cref{thm:normalization} identifies its value with the closure ratio.
\end{proof}
This is a finite-pivot theorem in exact arithmetic, not a polynomial bound
on the number of pivots, and not a guarantee for a floating-point variant
with a different pricing rule. Such variants require their own safeguards.
